\chapter{Numerabilidad}
\section{Axiomas de numerabilidad}
\subsection{Primer axioma de numerabilidad}
\begin{definition}[Primer axioma de numerabilidad]
Un espacio topológico $\displaystyle X $ se dice que es \textbf{primer axioma de numerabilidad}, y se escribe \textbf{I AN}, si todo punto $\displaystyle x \in X $ tiene una base de entornos $\displaystyle \mathcal{V}^{x} $ numerable.
\end{definition}
\begin{observation}
	Es fácil ver que si $\displaystyle \mathcal{V}^{x} = \left\{ V_{k}^{x}\right\} _{k \in \N} $, entonces $\displaystyle \left\{ U_{k}^{x}:=\Int\left(V_{k}^{x}\right)\right\} _{k\in \N} $ es una base numerable de entornos abiertos del punto $\displaystyle x $. 
\end{observation}
\begin{observation}
	Si $\displaystyle \mathcal{V}^{x}= \left\{ V_{k}^{x}\right\} _{k \in \N} $ es una base numerable de entornos de $\displaystyle x $, podemos formar una base numerable de entornos encajados
	\[\mathcal{W}^{x}= \left\{ W^{x}_{k}= V_{1}^{x}\cap \cdots \cap V^{x}_{k}\right\} _{k \in \N} .\]
	Claramente tenemos que $\displaystyle W^{x}_{k+1}\subset W^{x}_{k} $, $\displaystyle \forall k \in \N $. Análogamente, a partir de una base de entornos abiertos podemos encontrar una base numerable de entornos abiertos encajados. 
\end{observation}
\begin{eg}[Topologías metrizables]
	Todo espacio topológico $\displaystyle \left(X, \mathcal{T}\right) $ en el que la topología venga asociada a una distancia $\displaystyle d $ es I AN. En efecto, para cada $\displaystyle x \in X$ la colección $\displaystyle \left\{ B_{d}\left(x, \frac{1}{k}\right)\right\} _{k\in \N} $ es una base numerable de entornos abiertos de $\displaystyle x $. En particular, $\displaystyle \left(\R^{n}, \mathcal{T}_{u}\right) $ es I AN.
\end{eg}
\begin{eg}[Topología discreta y topología trivial]
	Es trivial que $\displaystyle \left(X, \mathcal{T}_{\text{dis}}\right) $ es I AN, puesto que para cada punto $\displaystyle x $ una base de entornos es $\displaystyle \left\{ x\right\}  $. Análogamente, $\displaystyle \left(X, \mathcal{T}_{\text{triv}}\right) $ también es I AN.
\end{eg}
\begin{eg}[Topología del punto]
	Dado un conjunto $\displaystyle X $ y $\displaystyle a \in X $, podemos considerar el espacio topológico $\displaystyle \left(X, \mathcal{T}_{a}\right) $. Claramente este espacio es I AN. En efecto, para $\displaystyle a $ podemos considerar la base de entornos $\displaystyle \left\{ a\right\}  $ y para $\displaystyle x \neq a $ basta con tomar $\displaystyle \left\{ a,x\right\}  $.
\end{eg}
\begin{eg}[Recta de Sorgenfrey]
La recta de Sorgenfrey $\displaystyle \left(\R, \mathcal{T}_{[,)}\right) $ también es I AN. En efecto, para $\displaystyle x \in \R $ basta con considerar la base de entornos $\displaystyle \left\{ \left[x, x + \frac{1}{n}\right)\right\} _{n \in \N} $.
\end{eg}
 
\begin{eg}[Topología que no es I AN]
Si $\displaystyle X $ es no numerable, la topología de los complementos finitos $\displaystyle \left(X, \mathcal{T}_{CF}\right) $ no es I AN. En efecto, un entorno abierto de $\displaystyle x\in X $ es de la forma $\displaystyle X \backslash A $ con $\displaystyle A \subset X $ finito y $\displaystyle x \not\in A $. 
Si tuviéramos una base numerable de entornos abiertos, la podríamos tomar encajada
\[\mathcal{W}^{x} = \left\{ W^{x}_{k} = X \backslash A_{k}\right\} _{k \in \N} ,\]
con los $\displaystyle A_{k} $ finitos. Pero entonces, 
\[\bigcap_{k \in \N}W_{k} = X \backslash \bigcup_{k \in \N}A_{k}\neq \emptyset .\]
Así, dado $\displaystyle x' \in \bigcap_{k \in \N}W_{k} $ con $\displaystyle x' \neq x $, el entorno abierto $\displaystyle X \backslash \left\{ x'\right\}  $ del punto $\displaystyle x $ no contiene a ninguno de los $\displaystyle W_{k} $, por lo que $\displaystyle \mathcal{W}^{x} $ no puede ser una base de entornos de $\displaystyle x $.
\end{eg}
\subsection{Sucesiones en espacios I AN}
Recordamos que una sucesión $\displaystyle \left\{ x_{k}\right\} _{k \in \N}\subset X $ la podemos ver como una aplicación
\[ \N \to X : k \to x_{k} .\]
\begin{definition}[Convergencia de una sucesión]
	La sucesión $\displaystyle \left\{ x_{k}\right\} _{k \in \N}\subset X $ converge a un punto $\displaystyle x \in X $, y escribimos $\displaystyle x_{k}\to x $, si para cada entorno abierto $\displaystyle U^{x} $ del punto $\displaystyle x $, existe un $\displaystyle k_{0}\in \N $ tal que $\displaystyle \forall k \geq k_{0} $, $\displaystyle x_{k}\in U^{x} $. 
\end{definition}
\begin{observation}
	Si $\displaystyle X $ es Hausdorff el límite de una sucesión, si existe, es único. En efecto, consideremos la sucesión $\displaystyle \left\{ x_{k}\right\} _{k \in \N} $ que converge a $\displaystyle x \in X $ y consideremos $\displaystyle x' \neq x $. Como $\displaystyle X $ es Hausdorff, existen $\displaystyle U^{x} $ y $\displaystyle U^{x'} $ entornos abiertos disjuntos. Como existe un $\displaystyle k_{0} \in \N $ tal que $\displaystyle \forall k \geq k_{0} $, $\displaystyle x_{k}\in U^{x} $, debe ser que $\displaystyle x_{k}\not\in U^{x'} $, $\displaystyle \forall k \geq k_{0} $, por lo que la sucesión no puede converger a $\displaystyle x' $. 
\end{observation}
\begin{prop}
Sea $\displaystyle \left(X, \mathcal{T}\right) $ un espacio topológico I AN y $\displaystyle A \subset X $. Las dos afirmaciones siguientes son equivalentes:
\begin{enumerate}
\item $\displaystyle x \in \overline{A} $.
\item Existe una sucesión $\displaystyle \left\{ x_{k}\right\} _{k\in \N}\subset A $ tal que $\displaystyle x_{k}\to x $.
\end{enumerate}
\end{prop}
\begin{proof}
La demostración es análoga a la que conocemos en espacios euclídeos.
\begin{description}
	\item[(1) $\displaystyle \Rightarrow $ (2)] Si $\displaystyle x \in \overline{A} $ tomamos una base de entornos abiertos encajados de $\displaystyle x $, $\displaystyle \left\{ U^{x}_{k}\right\} _{k\in \N} $. Para cada $\displaystyle k \in \N $ sabemos que $\displaystyle U^{x}_{k}\cap A\neq \emptyset $, y podemos escoger un elemento $\displaystyle x_{k}\in U^{x}_{k}\cap A $. Así, construimos una sucesión $\displaystyle \left\{ x_{k}\right\} _{k\in \N}\subset A $. 
	Para cada entorno abierto de $\displaystyle x $, $\displaystyle \Omega ^{x} $, existe $\displaystyle k_{0}\in \N $ tal que $\displaystyle U^{x}_{k}\subset \Omega ^{x} $, $\displaystyle \forall k \geq k_{0} $. Así, $\displaystyle x_{k}\in \Omega ^{x} $, $\displaystyle \forall k \geq k_{0} $ por lo que $\displaystyle x_{k}\to x $. 
\item[(2) $\displaystyle \Rightarrow $ (1)] Si hay una sucesión $\displaystyle \left\{ x_{k}\right\} _{k \in \N}\subset A $ que converge a $\displaystyle x $, entonces para cada entorno abierto de $\displaystyle x $, $\displaystyle \Omega ^{x} $, este contiene a todos los términos de la sucesión de uno en adelante. Como esos términos están en $\displaystyle A $ debe ser que $\displaystyle \Omega^{x}\cap A \neq \emptyset $.
\end{description}
\end{proof}
\subsection{Segundo axioma de numerabilidad}
\begin{definition}[Segundo axioma de numerabilidad]
Decimos que un espacio topológico es \textbf{segundo axioma de numerabilidad}, y lo escribimos \textbf{II AN}, si existe una base $\displaystyle \mathcal{B} $ numerable de abiertos de $\displaystyle X $.
\end{definition}
\begin{prop}
Si un espacio topológico $\displaystyle X $ es II AN, entonces es I AN.
\end{prop}
\begin{proof}
	Sea $\displaystyle \mathcal{B}= \left\{ B_{k}\right\} _{k\in \N } $ una base numerable de la topología. Para cada $\displaystyle x \in X $ el conjunto
	\[\mathcal{V}_{\mathcal{B}}^{x}= \left\{ B_{k}\in \mathcal{B} \; : \; x \in B_{k}\right\}  ,\]
	es una base de entornos abiertos de $\displaystyle x $. En efecto, si $\displaystyle \Omega ^{x} $ es entorno abierto de $\displaystyle x $, como todos los abiertos de $\displaystyle X $ debe ser unión de abiertos básicos. Como $\displaystyle x \in \Omega^{x} $ debe existir un $\displaystyle k\in \N $ tal que $\displaystyle x \in B_{k}\subset \Omega^{x} $.
\end{proof}
\begin{eg}[Falsedad del recíproco]
	El recíproco de la propiedad anterior no es cierto. En efecto, sea $\displaystyle X  $ un conjunto no numerable. Vimso que $\displaystyle \left(X, \mathcal{T}_{\text{dis}}\right) $ es I AN, pero no es II AN. En efecto, como todos los conjuntos $\displaystyle \left\{ x\right\}  $ son abiertos, una base de la topología debe incluirlos a todos. Como $\displaystyle X $ es no numerable, es imposible la existencia de una base numerable de $\displaystyle \mathcal{T}_{\text{dis}} $.
\end{eg}
\begin{eg}
El espacio topológico $\displaystyle \left(\R, \mathcal{T}_{u}\right) $ es II AN. En efecto, 
\[ \left\{ B\left(\left(q_{1}, \ldots, q_{n}\right), \frac{1}{k}\right) \; : \; \left(q_{1}, \ldots, q_{n}\right)\in \Q^{n}, \; k \in \N\right\}  ,\]
es una base numerable de $\displaystyle \mathcal{T}_{u} $. 
\end{eg}
\begin{eg}
	Si $\displaystyle X $ es no numerable y fijamos $\displaystyle a \in X $, tenemos que $\displaystyle \left(X, \mathcal{T}_{a}\right) $ no es II AN. En efecto, si $\displaystyle x \neq a $ tenemos que $\displaystyle \left\{ a,x\right\}  $ debe formar parte de cualquier base de la topología $\displaystyle \mathcal{T}_{a} $.
\end{eg}
\begin{eg}
La recta de Sorgenfrey no es II AN. En efecto, consideremos una base $\displaystyle \mathcal{B} $ para la topología de Sorgenfrey. Como para cada $\displaystyle x\in \R $ el intervalo $\displaystyle [x, x + 1) $ es un entorno abierto de $\displaystyle x $ en $\displaystyle \left(\R, \mathcal{T}_{[,)}\right) $, debe haber un abierto $\displaystyle B_{x}\in \mathcal{B} $ tal que $\displaystyle x \in B_{x}\subset[x,x +1) $. Observamos que $\displaystyle x = \inf B_{x} $. Para cualquier otro punto $\displaystyle x' \in \R $, procedemos igual y hallamos un abierto básico $\displaystyle B_{x'} $ con $\displaystyle x ' \in B_{x'}\subset [x',x'+1 )$. Como quiera que $\displaystyle x' = \inf B_{x'} $, necesariamente $\displaystyle B_{x} \neq B_{x'} $. Así, tenemos una aplicación inyectiva
\[ \R \to \mathcal{B} : x \to B_{x} ,\]
por lo que $\displaystyle \mathcal{B} $ no puede ser numerable. 
\end{eg}
\subsection{Separabilidad}
\begin{definition}[Espacio separable]
Un espacio topológico $\displaystyle X $ es \textbf{separable} si existe un subconjunto $\displaystyle A $ numerable denso.
\end{definition}
\begin{prop}
Si un espacio topológico $\displaystyle X $ es II AN, entonces es separable. 
\end{prop}
\begin{proof}
	Si $\displaystyle X $ es II AN y tomamos una base numerable de la topología $\displaystyle \mathcal{B} = \left\{ B_{k}\right\} _{k \in \N} $, podemos formar un conjunto $\displaystyle A = \left\{ a_{k}\right\} _{k \in \N} $ escogiendo un elemento $\displaystyle a_{k}\in B_{k} $ para cada $\displaystyle k \in \N $. Claramente, $\displaystyle A $ es numerable. Además es denso, puesto que dado un abierto no vacío $\displaystyle U $, como $\displaystyle U $ es unión de abiertos básicos debe ser que existe $\displaystyle k_{0}\in \N $ tal que $\displaystyle B_{k_{0}}\subset U $, por lo que $\displaystyle a_{k_{0}}\in A\cap U \neq \emptyset$.
\end{proof}
\begin{eg}
	Un espacio $\displaystyle X $ puede ser I AN y separable, pero no ser II AN. En efecto, si $\displaystyle X $ es un conjunto no numerable y tomando la topología $\displaystyle \mathcal{T}_{a} $, hemos visto que $\displaystyle \left(X, \mathcal{T}_{a}\right) $ es I AN. Además es numerable puesto que $\displaystyle \left\{ a\right\}  $ es denso y numerable. 
\end{eg}
\begin{eg}
Un espacio puede ser IAN y no separable. En efecto, si $\displaystyle X $ es no numerable basta con tomar $\displaystyle \left(X, \mathcal{T}_{\text{dis}}\right) $. En esta topología todos los subconjuntos son abiertos y cerrados, por lo que el único subconjunto denso es $\displaystyle X $.
\end{eg}
\begin{eg}
Un espacio puede ser separable y no ser I AN. Basta con considerar $\displaystyle \left(\R, \mathcal{T}_{CF}\right) $. Vimos que no era I AN, por otra parte tiene un subconjunto denso y numerable, por ejemplo $\displaystyle \Z $.
\end{eg}
\subsection{Espacios Lindelöf}
\begin{definition}[Recubrimiento y subrecubrimiento]
	Dada una colección de abiertos $\displaystyle \left\{ U_{i}\right\} _{i \in I}\subset X $ con $\displaystyle i \in I $, decimos que esa colección es un \textbf{recubrimiento por abiertos} de $\displaystyle X $ si $\displaystyle X = \bigcup_{i \in I}U_{i} $. Un \textbf{subrecubrimiento} es una subcolección $\displaystyle \left\{ U_{j} \; : \; j \in J \subset I\right\}  $ de modo que $\displaystyle X = \bigcup_{j \in J}U_{j} $.
\end{definition}
\begin{definition}[Espacio Lindelöf]
Un espacio topológico $\displaystyle X $ es de \textbf{Lindelöf} si cada recubrimiento por abiertos de $\displaystyle X $ admite un subrecubrimiento numerable.
\end{definition}
\begin{prop}
Si un espacio topológico $\displaystyle X $ es II AN, entonces es de Lindelöf. 
\end{prop}
\begin{proof}
	Si tenemos un recubrimiento por abiertos $\displaystyle X = \bigcup_{i \in I}U_{i} $, como sabemos que existe una base $\displaystyle \mathcal{B}= \left\{ B_{k}\right\} _{k\in \N} $ numerable, podemos proceder como sigue: para cada $\displaystyle B_{k} $, si está contenido en alguno de los abiertos del recubrimiento, escogemos uno de esos abiertos que lo contienen y lo denotamos por $\displaystyle U_{i\left(k\right)} $. Afirmamos que 
	\[X = \bigcup_{k \in \N, B_{k}\subset U_{i}\left(k\right)}U_{i\left(k\right)} .\]
	Si existiese 
	\[x_{0} \in X \backslash \bigcup_{k \in \N, B_{k}\subset U_{i\left(k\right)}}U_{i\left(k\right)} ,\]
	por ser la familia $\displaystyle \left\{ U_{i}\right\} _{i \in I} $ un recubrimiento por abiertos de $\displaystyle X $, debería existir un $\displaystyle U_{i\left(x_{0}\right)} $ de modo que $\displaystyle x_{0}\in U_{i\left(x_{0}\right)} $. Ahora bien, como $\displaystyle \mathcal{B}= \left\{ B_{k}\right\} _{k \in \N} $ es base de la topología, debería existir un $\displaystyle k_{0} $ tal que $\displaystyle x_{0}\in B_{k_{0}}\subset U_{i\left(x_{0}\right)} $. Pero entonces,
	\[x_{0}\in U_{i\left(k_{0}\right)}\subset \bigcup_{k \in \N, B_{k}\subset U_{i\left(k\right)}}U_{i\left(k\right)} ,\]
	que contradice nuestra suposición inicial.
\end{proof}
\begin{eg}
	Si $\displaystyle X $ es no numerable, $\displaystyle \left(X, \mathcal{T}_{a}\right) $ no es Lindelöf. En efecto, los conjuntos $\displaystyle U^{x}= \left\{ a,x\right\}  $ son abiertos. Si consideramos el recubrimiento 
	\[X = \bigcup_{x \in X \backslash \left\{ a\right\} }U^{x} ,\]
	no admite ningún subrecubrimiento.
\end{eg}
\section{Comportamiento por construcciones}
\begin{observation}
Es sencillo ver que todas las propiedades anteriores son propiedades topológicas, es decir, se conservan por homeomorfismos. 
\end{observation}

\begin{prop} Sea $\displaystyle X $ un espacio topológico. 
	\begin{enumerate}
	\item Si $\displaystyle X $ es I AN, todo subespacio $\displaystyle Y \subset X $ es I AN.
	\item Si $\displaystyle X $ es II AN, todo subespacio $\displaystyle Y \subset X $ es II AN. 
	\end{enumerate}
\end{prop}
\begin{proof}
	Sea $\displaystyle X $ un espacio topológico. 
	\begin{enumerate}
	\item Basta con observar que dado $\displaystyle y \in Y $, si $\displaystyle \mathcal{V}^{y}= \left\{ V^{y}_{k}\right\} _{k \in \N} $ es una base numerable de entornos abiertos de $\displaystyle y $ en $\displaystyle X $, entonces $\displaystyle \mathcal{W}^{y}= \left\{ V^{y}_{k}\cap Y\right\} _{k\in \N} $ es una base numerable de entornos de $\displaystyle y $ en $\displaystyle Y $.
	\item Dada una base numerable de la topología en $\displaystyle X $, $\displaystyle \mathcal{B}= \left\{ B_{k}\right\} _{k \in \N} $, sabemos que el conjunto $\displaystyle \mathcal{B}_{Y}= \left\{ B_{k}\cap Y\right\} _{k \in \N} $ es una base de la topología de $\displaystyle Y $ como subespacio.
	\end{enumerate}
\end{proof}
\begin{eg}
La recta de Sorgenfrey es de Lindelöf. Como cada abierto es unión de abiertos básicos, basta con considerar un recubrimiento de la forma
\[ \R = \bigcup_{i \in I}[a_{i}, b_{i}) ,\]
 donde cada uno de los $\displaystyle [a_{i}, b_{i}) $, con $\displaystyle a_{i} < b_{i} $, es un abierto de la base. Supongamos que consideramos la unión 
 \[\bigcup_{i \in I}\left(a_{i}, b_{i}\right) ,\]
 y sea $\displaystyle A $ el subconjunto de puntos de $\displaystyle \R $ que no están cubiertos por esa unión, es decir, $\displaystyle A = \R \backslash \bigcup_{i \in I}\left(a_{i}, b_{i}\right) $. 
 Si $\displaystyle x \in A $ debe existir $\displaystyle i \in I $ tal que $\displaystyle x = a_{i}\in [a_{i}, b_{i}) $. Sea $\displaystyle q_{x}\in \Q $ tal que $\displaystyle x < q_{x} < b_{i} $. Por otra parte, si $\displaystyle x' \neq x $ es otro punto de $\displaystyle A $, y $\displaystyle x < x' $, se tendrá de manera análoga que debe existir $\displaystyle j \in I $ tal que $\displaystyle x' = a_{j}\in [a_{j}, b_{j}) $. Además, podemos determinar un $\displaystyle q_{x'}\in \Q $ con $\displaystyle x' < q_{x'} < b_{j} $. 
 Se tiene que $\displaystyle q_{x} < q_{x'} $. En efecto, en caso contrario tendríamos que 
 \[x < x' < q_{x'} \leq q_{x} < b_{i} ,\]
 por lo que $\displaystyle x' \in \left(a_{i}, b_{i}\right) $, que contradice que $\displaystyle x ' \in A $. Así, hemos establecido una aplicación inyectiva
 \[A \to \Q : x \to q_{x} ,\]
 que muestra que $\displaystyle A $ es numerable. Si nos fijamos ahora en $\displaystyle \R\backslash A $, resulta que está recubierto por la unión de los intervalos abiertos $\displaystyle \left(a_{i}, b_{i}\right)\in I $, y pensando en $\displaystyle \R \backslash A $ como subespacio de la recta usual que es II AN, sabemos que nos podemos quedar con un subrecubrimiento numerable, los $\displaystyle \left\{ a_{j}, b_{j}\right\}  $ para $\displaystyle j \in J \subset I $ siendo $\displaystyle J $ numerable. 
 Hemos terminado, pues del recubrimiento inicial $\displaystyle \R = \bigcup_{i \in I}[a_{i}, b_{i}) $ hemos obtenido el subrecubrimiento numerable
 \[\R = \bigcup_{x \in A}[a_{i}, b_{i}) \cup \bigcup_{j \in J}[a_{j}, b_{j}) .\]
\end{eg}
\begin{prop}
Si $\displaystyle X $ es un espacio topológico separable y $\displaystyle Y \subset X $ abierto, entonces $\displaystyle Y $ es separable.
\end{prop}
\begin{proof}
Sea $\displaystyle A \subset X $ denso y numerable. Por densidad tenemos que $\displaystyle A \cap Y \neq \emptyset $. Además, el conjunto $\displaystyle A \cap Y $ es numerable y es denso en $\displaystyle Y $. En efecto, todo abierto en $\displaystyle Y $ lo es en $\displaystyle X $, por lo que $\displaystyle W \cap \left(A \cap Y\right)= W \cap A \neq \emptyset $.
\end{proof}
\begin{eg}
	Ser separable no es una propiedad hereditaria en general. Sean $\displaystyle \left(X, \mathcal{T}_{a}\right) $ con $\displaystyle X $ no numerable. Si $\displaystyle Y = X \backslash \left\{ a\right\}  $, entonces $\displaystyle Y $ hereda la topología discreta. Así, el único subconjunto denso es $\displaystyle Y $, que no es numerable. 
\end{eg}
\begin{prop}
Si el espacio topológico $\displaystyle \left(Z, \mathcal{T}\right) $ es Lindelöf y $\displaystyle C \subset X $ es cerrado, entonces $\displaystyle \left(C, \mathcal{T}|_{C}\right) $ es Lindelöf. 
\end{prop}
\begin{proof}
Si tomamos un recubrimiento por abiertos en $\displaystyle C $, 
\[C = \bigcup_{i \in I}V_{i}, \; V_{i} \in \mathcal{T}|_{C} .\]
Tenemos que existen $\displaystyle U_{i} \in \mathcal{T} $ tales que $\displaystyle U_{i} \cap C = V_{i} $, $\displaystyle \forall i \in I $. Entonces, $\displaystyle C \subset \bigcup_{i \in I}U_{i} $. Como $\displaystyle X \backslash C $ es abierto,
\[X = \left(X \backslash C\right)\cup \bigcup_{i \in I}U_{i} .\]
Como $\displaystyle X $ es Lindelöf, existe un subrecubrimiento numerable
\[X = \left(X \backslash C\right) \cup \bigcup_{j \in J}U_{j}, \]
con $\displaystyle J \subset I $ numerable. Claramente, $\displaystyle C = \bigcup_{j \in J}V_{j} $. 
\end{proof}
\begin{eg}[Lindelöf no es hereditaria]
	Sea $\displaystyle \left(Y, \mathcal{T}\right) $ un espacio topológico que sea $\displaystyle T_{1} $ pero no Lindelöf \footnote{Podría ser que $\displaystyle Y $ fuese no numerable y $\displaystyle \mathcal{T} =\mathcal{T}_{\text{dis}} $.} Consideramos $\displaystyle \omega \not\in Y $ y tomamos en $\displaystyle X = Y \cup \left\{ \omega \right\}  $ la topología $\displaystyle \tilde{\mathcal{T}} $ que tiene por abiertos los $\displaystyle U \subset X $ tales que o bien $\displaystyle U \in \mathcal{T} $ o bien $\displaystyle \omega \in U $ y $\displaystyle X \backslash U \subset Y $ es finito. \\
\indent Es obvio que de todo recubrimiento por abiertos de $\displaystyle X $ se puede encontrar uno numerable. En efecto, si 
\[ X = \bigcup_{i \in I}U_{i} ,\]
con $\displaystyle U_{i}\in \tilde{\mathcal{T}} $, $\displaystyle \forall i \in I $. Hay un $\displaystyle i_{0} \in I $ tal que $\displaystyle \omega \in U_{i_{0}} $, pero $\displaystyle X \backslash U_{i_{0}} $ es finito, por lo que además del $\displaystyle U_{i_{0}} $ para recubrir $\displaystyle X $ basta escoger $\displaystyle U_{i} $ en el que se encuentre cada punto de $\displaystyle X\backslash U_{i_{0}} $. \\
\indent Podemos observar que $\displaystyle Y $ con la topología que recibe como subespacio de $\displaystyle \left(X, \tilde{\mathcal{T}}\right) $ es el espacio $\displaystyle \left(Y, \mathcal{T}\right) $ de partida, que sabemos que no es Lindelöf. \\
\indent Para ver que $\displaystyle Y $ como subespacio no es de Lindelöf podemos tomar un recubrimiento de $\displaystyle Y $ por abiertos en $\displaystyle \mathcal{T} $,
\[Y = \bigcup_{\alpha \in \Lambda }V_{\alpha }, \; V_{\alpha }\in \mathcal{T} ,\]
que no admite subrecubrimiento numerable. Ahora bien, cada $\displaystyle V_{\alpha } $ es también abierto en $\displaystyle \left(X, \tilde{\mathcal{T}}\right) $. 
\end{eg}
\begin{eg}[I AN y II AN no es divisible]
Consideremos la relación de equivalencia en la que todos los enteros se identifican a un punto, y el correspondiente espacio cociente $\displaystyle \R / \Z $. Sea $\displaystyle p : \R \to \R/\Z $ la aplicación identificación de paso al cociente. Vamos a ver que $\displaystyle p\left(0\right) $ no tiene una base numerable de entornos abiertos, por lo que $\displaystyle \R / \Z $ no es I AN y tampoco es II AN. \\
\indent Si tomamos una familia numerable $\displaystyle \left\{ V_{k}^{p\left(0\right)}\right\} _{k \in \N} $ de entornos abiertos de $\displaystyle p\left(0\right) $, sabemos que para cada $\displaystyle k \in \N  $, $\displaystyle V_{k}^{p\left(0\right)}=p\left(U_{k}\right) $ donde $\displaystyle U_{k}\subset \R $ es un abierto usual de la recta $\displaystyle p $-saturado, es decir, con $\displaystyle \Z \subset U $. 
En particular, tenemos que existe un intervalo $\displaystyle \left(a_{k}, b_{k}\right)\subset U_{k} $ con 
\[k-1 < a_{k} < k < b_{k} < k+1 .\]
Tomemos un punto $\displaystyle x_{k}\neq k $ de forma que $\displaystyle x_{k}\in \left(a_{k}, b_{k}\right) $. Consideramos $\displaystyle \R \backslash \left\{ x_{k} \right\} _{k \in \N} $. Se trata de un conjunto abierto en la topología usual. Es $\displaystyle p $-saturado. Por tanto, $\displaystyle W^{p\left(0\right)}:= p\left(\R \backslash \left\{ x_{k}\right\} _{k \in \N}\right) $ es un entorno abierto de $\displaystyle p\left(0\right) $. 
Además, no contiene a ninguno de los $\displaystyle V^{p\left(0\right)}_{k} $, pues $\displaystyle p\left(x_{k}\right)\in V_{k}^{p\left(0\right)} $, pero $\displaystyle p\left(x_{k}\right)\not\in W^{p\left(0\right)} $. Así, es imposible que $\displaystyle p\left(0\right) $ tenga una base numerable de entornos abiertos. 
\end{eg}
\begin{prop}
Sea $\displaystyle f : X \to Y $ una aplicación continua, sobreyectiva y abierta.
\begin{enumerate}
\item Si $\displaystyle X $ es I AN, también lo es $\displaystyle Y $.
\item Si $\displaystyle X $ es II AN, también lo es $\displaystyle Y $.
\end{enumerate}
\end{prop}
\begin{proof}
	Supongamos que $\displaystyle f $ es una identificación abierta. 
	\begin{enumerate}
	\item Tomemos $\displaystyle y \in Y $. Por ser $\displaystyle f $ sobreyectiva tenemos que existe $\displaystyle x \in X $ tal que $\displaystyle f\left(x\right)=y $. Sea $\displaystyle \left\{ V_{k}^{x}\right\} _{k \in \N} $ una base numerable de entornos abiertos de $\displaystyle x $ en $\displaystyle X $. Como $\displaystyle f $ es abierta, tenemos que $\displaystyle \left\{ f\left(V_{k}^{x}\right)\right\} _{k\in \N} $ es una colección numerable de entornos abiertos de $\displaystyle y $. Veamos que es una base de entornos abiertos de $\displaystyle y $. 
		Si $\displaystyle \Omega^{y} $ es un entorno abierto de $\displaystyle y $, por continuidad de $\displaystyle f $ sabemos que $\displaystyle f^{-1}\left(\Omega ^{y}\right) $ es entorno abierto de $\displaystyle x $, por lo que existe $\displaystyle k_{0} \in \N $ tal que $\displaystyle V^{x}_{k_{0}}\subset f^{-1}\left(\Omega ^{y}\right) $, por lo que $\displaystyle f\left(V^{x}_{k_{0}}\right)\subset \Omega^{y} $.
	\item Sea $\displaystyle \mathcal{B}= \left\{ B_{k}\right\} _{k \in \N} $ una base numerable de la topología en $\displaystyle X $. Por ser $\displaystyle f $ abierta tenemos que $\displaystyle \mathcal{B}\left(f\right):= \left\{ f\left(B_{k}\right)\right\} _{k \in \N} $ es una colección numerable de abiertos en $\displaystyle Y $. Además, si $\displaystyle V \subset Y $ es abierto por la continuidad de $\displaystyle f $ se sigue que $\displaystyle f^{-1}\left(V\right)\subset X $ es abierto, por lo que se puede expresar como la unión de una subfamilia de los abiertos de $\displaystyle \mathcal{B} $. 
		Como $\displaystyle f $ es sobreyectiva se tiene que $\displaystyle f\left(f^{-1}\left(V\right)\right)=V $ y por tanto $\displaystyle V $ se expresa como unión de elementos de $\displaystyle \mathcal{B}\left(f\right) $. 
	\end{enumerate}
\end{proof}
\begin{prop}
Sea $\displaystyle f : X \to Y $ una aplicación continua y sobreyectiva. 
\begin{enumerate}
\item Si $\displaystyle X $ es separable, $\displaystyle Y $ también lo es. 
\item Si $\displaystyle X $ es Lindelöf, $\displaystyle Y $ también lo es.
\end{enumerate}
\end{prop}
\begin{proof}
Sea $\displaystyle f : X \to Y $ una aplicación continua y sobreyectiva. 
\begin{enumerate}
\item Sea $\displaystyle A \subset X $ numerable y denso. Se tiene que $\displaystyle f\left(A\right) $ es numerable. Por la continuidad de $\displaystyle f $, $\displaystyle f\left(\overline{A}\right)\subset \overline{f\left(A\right)} $ y como $\displaystyle \overline{A} = X $ y $\displaystyle f $ es sobreyectiva tenemos que $\displaystyle \overline{f\left(A\right)}= Y $, por lo que $\displaystyle f\left(A\right) $ es denso en $\displaystyle Y $.
\item Sea $\displaystyle Y = \bigcup_{i \in I}V_{i} $ un recubrimiento por abiertos de $\displaystyle Y $.
	Como $\displaystyle f $ es continua tenemos que $\displaystyle f^{-1}\left(V_{i}\right) $ es abierto en $\displaystyle X $. Además, $\displaystyle X = \bigcup_{i \in I}f^{-1}\left(V_{i}\right) $.
	Como $\displaystyle X $ es Lindelöf, existe $\displaystyle J \subset I $ numerable tal que $\displaystyle X = \bigcup_{j \in J}f^{-1}\left(V_{j}\right) $. Como $\displaystyle f $ es sobreyectiva se sigue que $\displaystyle Y = \bigcup_{j \in J}V_{j} $. 
\end{enumerate}
\end{proof}
\begin{prop}
Sean $\displaystyle X_{i} $, $\displaystyle i = 1, \ldots, r $ espacios topológicos. Las siguientes afirmaciones son equivalentes:
\begin{enumerate}
\item Los espacios $\displaystyle X_{i} $ son I AN para cada $\displaystyle i = 1, \ldots, r $. 
\item El producto $\displaystyle \prod^{r}_{i = 1}X_{i} $ es I AN. 
\end{enumerate}	
\end{prop}
\begin{proof}
Demostramos ambas implicaciones:
\begin{description}
	\item[(1) $\displaystyle \Rightarrow $ (2)] Si $\displaystyle \left(x_{1}, \ldots, x_{r}\right)\in \prod^{r}_{ i= 1}X_{i} $, podemos escoger bases numerables de entornos abiertos encajados $\displaystyle \mathcal{V}^{x_{i}}= \left\{ V_{k}^{x_{i}}\right\} _{k \in \N} $ para cada $\displaystyle x_{i} \in X_{i} $. Se puede comprobar que 
		\[ \left\{ W_{1} \times \cdots \times W_{r} \; : \; W_{i} \in \mathcal{V}^{x_{i}}\right\}  ,\]
		es una base numerable de entornos abiertos de $\displaystyle \left(x_{1}, \ldots, x_{r}\right) $. Es sencillo ver que es una colección numerable. Veamos que es base de entornos abiertos. Sea $\displaystyle \Omega^{\left(x_{1}, \ldots, x_{r}\right)} $ un entorno abierto de $\displaystyle \left(x_{1}, \ldots, x_{r}\right) $. Como las proyecciones son aplicaciones abiertas tenemos que $\displaystyle p_{i}\left(\Omega^{\left(x_{1}, \ldots, x_{r}\right)}\right) $ es entorno abierto de $\displaystyle x_{i} $ en $\displaystyle X_{i} $.
		Por tanto, existe $\displaystyle k_{i} \in \N $ tal que $\displaystyle V^{x_{i}}_{k}\subset p_{i}\left(\Omega^{\left(x_{1}, \ldots, x_{r}\right)}\right) $ para $\displaystyle k \geq k_{i} $. Claramente, 
		\[V^{x_{1}}_{k_{1}} \times \cdots \times V^{x_{r}}_{k_{r}}\subset \Omega^{\left(x_{1}, \ldots, x_{r}\right)} .\]
	\item[(2) $\displaystyle \Rightarrow $ (1)] Cogemos puntos $\displaystyle x_{i} \in X_{i} $. Tomemos una base numerable de entornos abiertos encajados de $\displaystyle \left(x_{1}, \ldots, x_{r}\right) $,
		\[\mathcal{W}^{\left(x_{1}, \ldots, x_{r}\right)} = \left\{ \Omega_{k}^{\left(x_{1}, \ldots, x_{r}\right)} \; \right\} _{k \in \N} .\]
		Por ser las proyecciones abiertas, para cada $\displaystyle i \in \left\{ 1, \ldots, r\right\}  $ se tiene que $\displaystyle V^{x_{i}}:=p_{i}\left(\Omega_{k}^{\left(x_{1}, \ldots, x_{r}\right)}\right) $ es un entorno abierto de $\displaystyle x_{i} $ en $\displaystyle X_{i} $. Es más, 
		\[ \left\{ p_{i}\left(\Omega_{k}^{\left(x_{1}, \ldots, x_{r}\right)}\right)\right\} _{k \in \N} ,\]
	es una base encajada de abiertos de $\displaystyle x_{i} $. Es sencillo ver que están encajados. Veamos que efectivamente es base de entornos: si $\displaystyle U^{x_{i}} $ es entorno abierto de $\displaystyle x_{i} $ se tiene que 
	\[p^{-1}_{i}\left(U^{x_{i}}\right)= X_{1} \times \cdots \times X_{i-1} \times U^{x_{i}} \times X_{i+1} \times \cdots \times X_{r} ,\]
	es abierto y es entorno de $\displaystyle \left(x_{1}, \ldots, x_{r}\right) $ por lo que existe $\displaystyle k \in \N $ tal que $\displaystyle \Omega_{k}^{\left(x_{1}, \ldots, x_{r}\right)}\subset p^{-1}_{i}\left(U^{x_{i}}\right) $. Proyectando sobre el factor $\displaystyle i $-ésimo obtenemos que $\displaystyle p_{i}\left(\Omega^{\left(x_{1}, \ldots, x_{r}\right)}_{k}\right)\subset U^{x_{i}} $.
\end{description}
\end{proof}
\begin{prop}
Sean $\displaystyle X_{i} $ espacios topológicos para $\displaystyle i = 1, \ldots, r $. Son equivalentes:
\begin{enumerate}
\item Los espacios $\displaystyle X_{i} $ son II AN para cada $\displaystyle i = 1, \ldots, r $.
\item El producto $\displaystyle \prod^{r}_{i = 1}X_{i} $ es II AN.
\end{enumerate}
\end{prop}
\begin{proof}
Consideramos ambas implicaciones.
\begin{description}
	\item[(1) $\displaystyle \Rightarrow $ (2)] Si tomamos bases numerables $\displaystyle \mathcal{B}_{i} = \left\{ B^{i}_{k}\right\} _{k \in \N} $ para las topologías en $\displaystyle X_{i} $, sabemos que 
		\[\mathcal{B} = \left\{ B^{1}_{k_{1}} \times \cdots \times B^{r}_{k_{r}} \right\} _{k_{1}, \ldots, k_{r} \in \N} ,\]
		es una base de la topología producto y claramente es numerable. 
	\item[(2) $\displaystyle \Rightarrow $ (1)] Sea $\displaystyle \mathcal{B} $ es una base numerable de la topología producto con la forma
		\[\mathcal{B}= \left\{ W_{k}\right\} _{k\in \N} .\]
		Por ser las proyecciones abiertas tenemos que $\displaystyle p_{i}\left(W_{k}\right) $ es abierto en $\displaystyle X_{i} $, $\displaystyle \forall k \in \N $. Además, $\displaystyle \left\{ p_{i}\left(W_{k}\right)\right\} _{k \in \N} $ es base de la topología en $\displaystyle X_{i} $. En efecto, si $\displaystyle \Omega_{i} $ es abierto en $\displaystyle X_{i} $ tenemos que $\displaystyle p^{-1}_{i}\left(\Omega_{i}\right) $ es abierto en $\displaystyle \prod^{r}_{ i= 1}X_{i} $. 
		Por tanto, podemos escribir
		\[p^{-1}_{i}\left(\Omega_{i}\right)=\bigcup_{n\in \Lambda \subset \N}W_{n} .\]
		Tenemos claramente que 
		\[\Omega_{i}=\bigcup_{n \in \Lambda \subset \N}p_{i}\left(W_{n}\right) .\]
\end{description}
\end{proof}
\begin{prop}
Sean $\displaystyle X_{i} $ espacios topológicos con $\displaystyle i = 1, \ldots, r $. Son equivalentes:
\begin{enumerate}
\item Los espacios $\displaystyle X_{i} $ son separables para cada $\displaystyle i = 1, \ldots, r $. 
\item El producto $\displaystyle \prod^{r}_{i = 1}X_{i} $ es separable. 
\end{enumerate}
\end{prop}
\begin{proof}
Consideramos ambas implicaciones.
\begin{description}
\item[(1) $\displaystyle \Rightarrow $ (2)] Si para cada $\displaystyle i = 1, \ldots, r $ existe un conjunto $\displaystyle A_{i}\subset X_{i} $ que es denso y numerable, entonces claramente tendremos que $\displaystyle A_{1} \times A_{r} $ es numerable y denso en $\displaystyle \prod^{r}_{i = 1}X_{i} $ puesto que 
	\[ \overline{A_{1} \times \cdots \times A_{r}}= \overline{A}_{1} \times \cdots \times \overline{A}_{r} = X_{1} \times \cdots \times X_{r} = \prod^{r}_{i = 1}X_{i} .\]
\item[(2) $\displaystyle \Rightarrow $ (1)] Si $\displaystyle A \subset \prod^{r}_{i= 1}X_{i} $ es numerable y denso, para cada $\displaystyle i \in \left\{ 1, \ldots, r\right\}  $ tenemos que $\displaystyle p_{i}\left(A\right) $ es numerable. Tomemos $\displaystyle U_{i}\subset X_{i} $ un abierto no vacío. Tenemos que $\displaystyle p^{-1}_{i}\left(U_{i}\right) $ es un abierto no vacío de $\displaystyle \prod^{r}_{i = 1}X_{i} $. Por tanto $\displaystyle p^{-1}_{i}\left(U_{i}\right)\cap A \neq \emptyset $. Si $\displaystyle a = \left(a_{1}, \ldots, a_{r}\right)\in p^{-1}_{i}\left(U_{i}\right)\cap A $, entonces $\displaystyle a_{i} \in U_{i} \cap p_{i}\left(A\right) $. 
\end{description}
\end{proof}
\begin{eg}[Lindelöf no es productiva]
	Consideramos la recta de Sorgenfrey $\displaystyle \left(\R, \mathcal{T}_{[,)}\right) $, que sabemos que es Lindelöf. Veamos que $\displaystyle \left(\R, \mathcal{T}_{[,)}\right) \times \left(\R, \mathcal{T}_{[,)}\right) $ no es Lindelöf. Si lo fuera el subconjunto cerrado $\displaystyle y = \left\{ x + y = 0\right\}  $ sería Lindelöf. Sin embargo, hereda la topología discreta y es no numerable por lo que no es Lindelöf.
\end{eg}
\begin{prop}
Sean $\displaystyle X_{i} $ espacios topológicos con $\displaystyle i = 1, \ldots, r $ y consideremos el espacio suma $\displaystyle \sum^{r}_{i = 1}X_{i} $. 
\begin{enumerate}
	\item $\displaystyle X_{i} $ es I AN $\displaystyle \forall i \in \left\{ 1, \ldots, r\right\}  $ $\displaystyle \iff  $ $\displaystyle \sum^{r}_{i = 1}X_{i} $ es I AN. 
	\item $\displaystyle X_{i} $ es II AN $\displaystyle \forall i \in \left\{ 1, \ldots, r\right\}  $ $\displaystyle \iff  $ $\displaystyle \sum^{r}_{i = 1}X_{i} $ es II AN. 
	\item $\displaystyle X_{i} $ es separable $\displaystyle \forall i \in \left\{ 1, \ldots, r\right\}  $ $\displaystyle \iff  $ $\displaystyle \sum^{r}_{i = 1}X_{i} $ es separable.
	\item $\displaystyle X_{i} $ es Lindelöf $\displaystyle \forall i \in \left\{ 1, \ldots, r\right\}  $ $\displaystyle \iff  $ $\displaystyle \sum^{r}_{i = 1}X_{i} $ es Lindelöf. 
\end{enumerate}
\end{prop}

